\section{Limitations and Model Assumptions}
Our current implementation using Cosserat rod theory captures the full 3D kinematics (twist, shear) but relies on several simplifying assumptions. Primarily, we modeled the spine as a continuous, structurally isotropic or uniformly anisotropic rod. In reality, the spine exhibits highly complex, localized anisotropy. Our current implementation treats protein dynamics with simplified kinetics and does not fully capture mechanochemical feedback. Direct experimental measurement of mechanosensory protein turnover rates during growth spurts would refine predictions. Additionally, our tensor coupling model assumes orthotropic symmetry; future work should incorporate full anisotropic elasticity tensors measured via multi-directional mechanical testing.

Additionally, our protein dynamics model treats the complex metabolic network as a simplified two-component system (Supply vs. Demand). While this captures the essential scaling behavior, it abstracts away specific pathway kinetics (e.g., HIF-1$\alpha$ switching, NAD+ depletion). Finally, the mapping from discrete genetic expression to the continuous $I(s)$ field remains phenomenological; future refinements will link this directly to single-cell transcriptomic atlases of the developing spine.
